% Generated draft: requires mathematical and source review.
\documentclass[11pt]{article}
\usepackage[margin=25mm]{geometry}
\usepackage{fontspec}
\setmainfont{Latin Modern Roman}
\usepackage{amsmath,amssymb,mathtools,unicode-math}
\setmathfont{Latin Modern Math}
\usepackage{enumitem,needspace,xurl,hyperref}
\setlength{\parindent}{0pt}
\setlength{\parskip}{5pt}
\setlength{\emergencystretch}{3em}
\newcommand{\N}{\mathbb{N}}
\newcommand{\Z}{\mathbb{Z}}
\newcommand{\Q}{\mathbb{Q}}
\newcommand{\R}{\mathbb{R}}
\newcommand{\C}{\mathbb{C}}
\newcommand{\F}{\mathbb{F}}
\newcommand{\A}{\mathbb{A}}
\newcommand{\PP}{\mathbb{P}}
\newcommand{\E}{\mathbb{E}}
\newcommand{\eps}{\varepsilon}
\newcommand{\dd}{\,\mathrm{d}}
\begin{document}
\section*{UnsolvedMath ID: 30001546}
\textbf{Source problem code:} catalogue entry
\subsection{Definitions and mathematical statement}
\textbf{Supplied question.} ``How can sequences with large variations or fluctuations be detected and defined?''

\textbf{Formulation limits.} The supplied record does not specify the index set, the space of sequence values, the class of admissible sequences, or any mathematical meaning of ``variation,'' ``fluctuation,'' or ``large.'' It also does not specify whether ``detection'' means a mathematical characterization, an algorithm, or a statistical procedure, nor what information about a sequence is available for detection. Consequently, no precise predicate, boundary conditions, or quantified mathematical assertion can be extracted without additional hypotheses. The question is retained as a request to formulate suitable definitions and detection criteria; no particular definition or criterion is imposed here.

The registry labels this question open. Its supplied literature assessment describes it as a broad question rather than a precise theorem and reports no general resolution in the cited literature. These are supplied curation claims, not independently verified mathematical facts.
\subsection{Short English statement}
The question asks how to define and recognize sequences exhibiting large fluctuations. The sequence domain, the measure of fluctuation, the meaning of largeness, and the requirements for detection remain unspecified, so a precise mathematical problem requires further source clarification.
\subsection{Sources}
\begin{itemize}
\item Local UnsolvedMath registry, permanent ID \(30001546\), source problem code \texttt{catalogue entry}, title ``Detection of Highly Fluctuating Sequences,'' in the Oberwolfach Reports Open Problems collection. The supplied proposed year is 2010; no proposer is identified.
\item Supplied report citation: \emph{Mini-Workshop: Combinatorics on Words (2011)}, Oberwolfach Reports, DOI 10.4171/owr/2010/37, pp.\ 2195--2244. Supplied source URL and literature-source link, labeled ``original OWR report'': \url{https://doi.org/10.4171/owr/2010/37}. The background describes Oberwolfach Reports as documenting mathematical workshops and their open-problem discussions. The report is supplied as provenance; no passage establishing the intended definitions or detection task is included.
\item Literature assessment checked on \texttt{2026-08-21}: supplied status ``open.'' The supplied assessment states that this broad question about detecting sequences with large fluctuations is not a precise theorem, has no general resolution in the cited literature, and remains open with the report as provenance. This is supplied dated curation, not an independently verified literature assessment; its time-sensitive status and cited claims require independent verification.
\item The registry marks the statement as ``exact,'' with the supplied verification note that it is internally readable and was retained byte-for-byte. This does not establish verification against the cited report.
\end{itemize}
\end{document}
